\documentclass[10pt,a4paper]{amsart}
\usepackage[margin=19mm]{geometry}
\usepackage{amsmath,amssymb,mathtools}
\usepackage[scr=rsfs,cal=euler]{mathalfa}
\usepackage{xfrac}
\usepackage[unicode,hidelinks,bookmarks=false]{hyperref}
\newcommand{\ie}{i.e.\ }
\newcommand{\cf}{cf.\ }
\newcommand{\resp}{respectively\ }
\newcommand{\Subsection}{\subsection}
\providecommand{\nobreakdash}{\nobreak\hbox{-}}
\setlength{\emergencystretch}{2em}
\multlinegap0pt
\usepackage{amssymb}
\usepackage{mathtools}
\newtheorem{lemma}{Lemma}
\newcommand{\R}{\mathbb R}
\newcommand{\T}{\mathbb T}
\title{Analyticity of Lyapunov Exponents for Mixed Markov--Quasi-Periodic Cocycles}
\author{EL Hadji Yaya Tall}
\date{Source: arXiv:2608.01569v1 (CC BY 4.0)}
\begin{document}
\maketitle

\noindent\emph{Source and licence.} Analyticity of Lyapunov Exponents for Mixed Markov--Quasi-Periodic Cocycles. Version arXiv:2608.01569v1 (CC BY 4.0), distributed by arXiv under the Creative Commons Attribution 4.0 International licence, http://creativecommons.org/licenses/by/4.0/, as recorded in the arXiv licence metadata for this exact version and shown on the abstract page https://arxiv.org/abs/2608.01569v1. Full licence notice: \texttt{licenses/arxiv-2608.01569-CC-BY-4.0.txt}.\par\smallskip

\noindent\textit{Editorial scope.} Counted unit: the article's lemma lem:operator-bounded (source0.tex lines 2494--2529). The article's complete original proof of it is delivered with it (source0.tex lines 2531--2678, one \texttt{proof} environment of 148 lines). No mathematics is written, repaired or re-derived: the statement and the proof are the article's own lines, in the article's own notation, and no retained line is substituted or re-flowed.  Retained uncounted as the source-local context the proof needs: lines 2423--2432.  Nothing was omitted from any retained span, so the package carries no omission record.  No retained line contains a citation, so the wrapper carries no bibliography and no bibliography entry.  No retained line contains a figure environment, an image-inclusion command or drawing code, so no image is delivered and the package declares no asset dependency.  Editorial formatting only, in addition: an AMS \texttt{amsart} preamble that restates the article's own declarations and macros used by the retained lines, namely the environments \texttt{lemma} on separate counters, together with the standard packages those lines need.  The retained environments print in this document as Lemma 1 in order of appearance, whereas the article prints the same items with the numbering of its own full document.  The complete native source of the article as distributed in the arXiv e-print for this exact version is bundled unchanged as \texttt{sources/206641/source0.tex}.
\begin{equation}\label{eq:forward-Markov-operator}
 (\mathcal T_Z\Phi)(i,t, v)
 =
 \sum_{j=1}^Nz_{ij}\,
 \Phi\left(
 j,\,
 t+\theta_i,\,
 A_i(t)v
 \right),
\end{equation}

\begin{lemma}\label{lem:operator-bounded}
For every \(Z\in E(U)\), the operator
\[
 \mathcal T_Z:\mathcal B_\alpha\longrightarrow\mathcal B_\alpha
\]
is bounded.  More precisely, if
\[
 M(Z)=\max_{i\in X}\sum_{j=1}^N|z_{ij}|
\]
%and
%\begin{equation}\label{eq:uniform-condition-number}
% \kappa
% =
% \max_{\substack{i\in X\\t\in\T^m}}
% \|A_i(t)\|\,\|A_i(t)^{-1}\|,
%\end{equation}
then
\begin{equation}\label{eq:operator-norm-bound}
 \|\mathcal T_Z\|_{\mathcal L(\mathcal B_\alpha)}
 \leq
 M(Z)\max\{1,\kappa^{2\alpha}\}.
\end{equation}
Consequently, if \(V\subset E(U)\) is bounded, then
\[
 \sup_{Z\in V}
 \|\mathcal T_Z\|_{\mathcal L(\mathcal B_\alpha)}
 <\infty.
\]
Moreover,
\[
 E(U)\ni Z\longmapsto
 \mathcal T_Z\in\mathcal L(\mathcal B_\alpha)
\]
is complex-linear and hence entire.  In particular, its restriction
to \(H(U)\) is holomorphic.
\end{lemma}

\begin{proof}
First observe that \(\mathcal T_Z\Phi\) is continuous.  This follows
from the continuity of \(\Phi\), the continuity of \(t\mapsto A_i(t)\),
and the fact that the sum in
\eqref{eq:forward-Markov-operator} is finite.

For the uniform norm, we have
\begin{align}
 |(\mathcal T_Z\Phi)(i,t, v)|
 &\leq
 \sum_{j=1}^N|z_{ij}|
 \left|
 \Phi\left(
 j,t+\theta_i,A_i(t)v
 \right)
 \right|
 \notag\\
 &\leq
 \left(\sum_{j=1}^N|z_{ij}|\right)
 \|\Phi\|_\infty.
\end{align}
Taking the supremum over \(i,t, v\) gives
\begin{equation}\label{eq:operator-sup-bound}
 \|\mathcal T_Z\Phi\|_\infty
 \leq
 M(Z)\|\Phi\|_\infty.
\end{equation}

We next estimate the projective Hölder seminorm.%  Recall that, for
%nonzero representatives \(u,v\in\R^d\),
%\[
% d(\bar u,\bar v)
% =
% \frac{\|u\wedge v\|}{\|u\|\,\|v\|}.
%\]
%For every \(A\in\operatorname{GL}_d(\R)\),
%\begin{align*}
% d(\overline{Au},\overline{Av})
% &=
% \frac{\|Au\wedge Av\|}{\|Au\|\,\|Av\|}\\
% &\leq
% \frac{
% \|\bigwedge^2A\|\,\|u\wedge v\|
% }{
% \|A^{-1}\|^{-2}\|u\|\,\|v\|
% }\\
% &\leq
% \|A\|^2\|A^{-1}\|^2
% d(\bar u,\bar v).
%\end{align*}
%Applying this with \(A=A_i(t)\) yields
%\begin{equation}\label{eq:projective-Lipschitz-bound}
% d\left(
% \overline{A_i(t)u},
% \overline{A_i(t)v}
% \right)
% \leq
% \kappa^2d(\bar u,\bar v).
%\end{equation}
Let \( u\neq  v\).  We have that
%\eqref{eq:projective-Lipschitz-bound}, we obtain
\begin{align*}
 &\left|
 (\mathcal T_Z\Phi)(i,t,u)
 -
 (\mathcal T_Z\Phi)(i,t, v)
 \right|
 \\
 &\quad\leq
 \sum_{j=1}^N|z_{ij}|
 \left|
 \Phi\left(
 j,t+\theta_i, A_i(t)u
 \right)
 -
 \Phi\left(
 j,t+\theta_i,A_i(t)v
 \right)
 \right|
 \\
 &\quad\leq
 \left(\sum_{j=1}^N|z_{ij}|\right)
 [\Phi]_\alpha
 d\left( A_i(t)u, A_i(t)v
 \right)^\alpha
 \\
 &\quad\leq
 M(Z)\kappa^{2\alpha}
 [\Phi]_\alpha
 d(u, v)^\alpha.
\end{align*}
It follows that
\begin{equation}\label{eq:operator-holder-bound}
 [\mathcal T_Z\Phi]_\alpha
 \leq
 M(Z)\kappa^{2\alpha}[\Phi]_\alpha.
\end{equation}
Combining \eqref{eq:operator-sup-bound} and
\eqref{eq:operator-holder-bound}, we get
\begin{align*}
 \|\mathcal T_Z\Phi\|_{\mathcal B_\alpha}
 &\leq
 M(Z)\|\Phi\|_\infty
 +
 M(Z)\kappa^{2\alpha}[\Phi]_\alpha\\
 &\leq
 M(Z)\max\{1,\kappa^{2\alpha}\}
 \|\Phi\|_{\mathcal B_\alpha}.
\end{align*}
This proves \eqref{eq:operator-norm-bound}.

If \(V\subset E(U)\) is bounded, then, because \(E(U)\) is finite
dimensional,
\[
 M_V:=\sup_{Z\in V}M(Z)<\infty.
\]
Hence
\[
 \sup_{Z\in V}
 \|\mathcal T_Z\|_{\mathcal L(\mathcal B_\alpha)}
 \leq
 M_V\max\{1,\kappa^{2\alpha}\}<\infty.
\]

Finally, for each admissible edge \((i,j)\in\mathcal E\), define the
bounded operator \(\mathcal L_{ij}\) by
\[
 (\mathcal L_{ij}\Phi)(r,t, v)
 =
 \mathbf 1_{\{i\}}(r)\,
 \Phi\left(
 j,t+\theta_i,A_i(t)v \right).
\]
Then
\begin{equation}\label{eq:operator-linear-decomposition}
 \mathcal T_Z
 =
 \sum_{(i,j)\in\mathcal E}z_{ij}\mathcal L_{ij}.
\end{equation}
This is a finite sum of bounded operators and shows that
\[
 Z\longmapsto\mathcal T_Z
\]
is complex-linear from \(E(U)\) into
\(\mathcal L(\mathcal B_\alpha)\).  In particular, it is entire in
the operator norm.  Its restriction to the affine subspace \(H(U)\)
is therefore holomorphic.
\end{proof}

\end{document}
